\section{Limitations and Conclusion}

The current evidence is deliberately scoped in three ways. First, the validity and calibration scope is empirical: \fastcp{} is post-hoc, and local weighting with augmented calibration pooling does not restore finite-sample conformal guarantees under arbitrary distribution shift. The expanded benchmark therefore supports a coverage-efficiency tradeoff, not a free calibration improvement: prediction sets are significantly smaller, but mean coverage and undercoverage-tail behavior are worse than augmented split conformal prediction.

Second, the data and model coverage remains bounded. The strongest results are at $\alpha=0.05$ on 35 small-to-medium univariate UCR datasets with synthetic corruptions. We add $\alpha\in\{0.1,0.2\}$, held-out $\lambda$ Pareto analysis, UEA multivariate validation, MiniROCKET, and FCN checks, but these supplements do not replace real corrupted deployment datasets or a full deep TSC study over InceptionTime, ResNet-style models, and transformers.

Third, the fingerprint and hyperparameter design is heuristic. Confidence and margin features are useful but may introduce circularity, because the classifier's own uncertainty helps determine the threshold. The bandwidth $h$ and mix $\lambda$ have no adaptive error-control rule; under independent coverage-prioritized validation, the selected value is $\lambda=1$, so intermediate lambdas should be reported as Pareto tradeoff points rather than validated defaults. The random-fingerprint and feature-ablation controls further show that the full hand-designed fingerprint is not uniquely responsible for the efficiency gain: confidence-only weighting is safer, input-only weighting is more aggressive, and random weighting still contributes some set-size reduction.

Despite these limits, the results support a practical conclusion. Perturbation-aware local weighting, combined with global fallback calibration, provides a lightweight way to adapt conformal thresholds to corrupted time-series inputs. Coverage-prioritized validation preserves the augmented global threshold when coverage is the primary objective; the Pareto-oriented setting improves prediction-set efficiency with a measurable coverage cost. For practitioners working under modest compute budgets, \fastcp{} offers a reproducible calibration layer when this tradeoff is acceptable.
